% What this holds: \part{Historical Background} and its lead-in paragraph.
% Read by arlington-bsap.tex. Prose lives here; formatting lives in the wrapper and style/.
\part{Historical Background}

% Lead-in to Part A: why the history begins in 1870, and the county the Board
% inherited. The footnote names who governed the land before
% (docs/questions.csv, before-1870-governance).
The Arlington County Board's history begins in 1870, and so does this report's.
Virginia's constitution of 1869 created a board of supervisors in every county,
and Alexandria County, as Arlington was called until 1920, elected its first
three in May
1870.\autocite[art.~VII, sec.~2]{vaconstitution1869}\fnsep\autocite[chs.~39, 76 and 188]{vaacts1870}\fnsep\autocite[note~2]{forstall1996}
The county itself was then 69 years old: Fairfax County's land until 1801, then
the Virginia side of the federal district, which Congress governed through a
circuit court and appointed magistrates until it returned the county to
Virginia in
1847.\autocite{usstat1790residence}\fnsep\autocite[sec.~2]{usstat1801organic}\fnsep\autocite[secs.~1, 4 and 7]{usstat1801supplementary}\fnsep\autocite{usstat1846retrocession}
The Board is the government this report is about, not the first to have ruled
this ground.\footnote{The Nacotchtank, whom the English also called Necostins,
lived on this shore of the Potomac when John Smith charted the river in 1608;
his map places a village, Nameroughquena, in what is now Arlington, and by 1679
they had left it, to disease, war and the loss of their fishing, hunting and
farming grounds (\cite{arlingtonva2011markers}; \cite{nps2025rooseveltisland}).
The Virginia colony and then Fairfax County governed the land until Congress
accepted a federal district on the Potomac in 1790, and the district's Virginia
side was a county of the District of Columbia from 1801 until Congress returned
it to Virginia.}
